\section{EvoDiff：进化尺度的蛋白质序列扩散模型}

\subsection{模型概述与训练数据}

EvoDiff 是 Amini 团队在 MSR 开发的\textbf{基于离散扩散的蛋白质生成模型}，专门为功能性蛋白质设计而构建。该项目由 Sarah Alamdari 主导开发，Amini 与同事 Kevin Yang 共同指导研究方向。

\begin{importantbox}{EvoDiff 的核心特征}
\begin{itemize}
    \item \textbf{模型类型}：基于离散扩散（Discrete Diffusion）的生成模型
    \item \textbf{数据模态}：蛋白质氨基酸序列（离散 token）
    \item \textbf{训练数据}：约 5000 万条独特蛋白质序列，跨越进化尺度
    \item \textbf{核心能力}：从全 mask 状态逐步去噪生成新的蛋白质序列
    \item \textbf{扩展能力}：支持功能条件生成（motif scaffolding）
\end{itemize}
\end{importantbox}

"进化尺度"意味着训练数据涵盖了生命树 (tree of life) 上不同生物体的蛋白质序列——从细菌到人类，包含了极其多样化的生物功能表征。这种广泛的数据覆盖使模型能够学习自然蛋白质的普遍分布规律。

\subsection{两种序列输入模式}

EvoDiff 支持两种序列输入模式，分别适用于不同的建模场景：

\textbf{模式一：单序列 (Single Sequence)}

直接对单条蛋白质序列进行离散扩散建模。适用于无条件的蛋白质从头设计 (de novo design)。

\textbf{模式二：Multiple Sequence Alignment (MSA)}

利用进化上下文信息。对于给定的目标序列，通过序列相似性搜索找到一组\textbf{相关但不完全相同的进化同源序列}，将这些序列对齐排列成一个矩阵（MSA），然后在整个 MSA 矩阵上执行离散扩散。

\begin{knowledgebox}{MSA 为什么能提供额外信息？}
Multiple Sequence Alignment 是计算生物学中的经典概念。进化上相关的蛋白质序列虽然在某些位置发生了突变，但在功能关键的位置上往往高度保守。MSA 中的列级保守模式 (conservation pattern) 隐含地编码了\textbf{哪些位置对功能至关重要}这一关键信息。EvoDiff 的 MSA 模式通过联合建模多条相关序列，能够从进化信号中提取更丰富的先验知识，指导生成过程。
\end{knowledgebox}

\subsection{生成过程的可视化}

EvoDiff 的生成过程可以直观地理解为：

\begin{enumerate}
    \item \textbf{初始化}：全 \texttt{[MASK]} 序列，长度为目标蛋白质的长度
    \item \textbf{迭代去噪}：每一步，模型预测并填充一部分 mask 位置的氨基酸
    \item \textbf{序列成形}：随着步骤推进，越来越多的位置被填充，序列逐渐成形
    \item \textbf{完成}：所有位置都被填充，得到一条完整的蛋白质序列
\end{enumerate}

Amini 在讲座中展示了一个动态可视化：左侧是 EvoDiff 序列模型的逐步生成过程，右侧同步展示了当前序列对应的预测三维结构。值得强调的是，\textbf{EvoDiff 的学习和生成完全在序列空间进行}，没有利用任何结构信息——右侧的结构仅用于最终可视化。

\subsection{本章小结}

EvoDiff 利用离散扩散框架，在约 5000 万条进化尺度的蛋白质序列上进行训练，支持单序列和 MSA 两种输入模式。模型完全在序列空间中操作，从全 mask 状态逐步去噪生成新蛋白质。MSA 模式通过引入进化上下文信息，为生成过程提供更丰富的先验知识。


